

% Positioning. Background.
Sustainable AI has emerged in an era where AI workload enormously increase the data center emissions~\cite{jha-etal:2025}. 
Computationally expensive and time-consuming neural networks are being trained daily with massive datasets, while the demand for even larger, better or more task-oriented datasets shows no sign of slowing down.
Questiong like: "How much more data do I need~\cite{mahmood-etal:2022}?" and "do we need more training data~\cite{zhu-vondrick-etal:2016}?" highlight that large-scale datasets have increasingly become a contraint rathen than a beneficial resource. 


In the context of dataset reduction, two main research directions have emerged:
\emph{coreset selection} and \emph{dataset distillation}. 
Coreset methods work directly on real samples and aim to pick a small subset
that preserves the information characteristics of the full dataset, typically by
optimizing geometry, or gradient similarity.
On the other hand, distillation~\cite{lei-tao:2024} synthesizes a tiny set of artificial dataset whose training
dynamics mimic those of the original dataset, often using generative models or 
trajectory matching in pixel or latent space.

Recent work~\cite{moser-shanbhag-etal:2025} has highlighted several structural limitations of existing coreset methods.
Training-free, geometry-based approaches such as Herding~\cite{welling:2009} or 
$k$-Center Greedy~\cite{sener-savarese:2017} crucially depend on the quality of the feature extractor. 
If the embedding does not faithfully capture the class structure, selection tends to pick arbitrary
points with good geometric coverage but poor alignment with the final decision boundary.
Training-oriented approaches such as CAL~\cite{margatina-vernikos-barrault:2021} and 
DeepFool~\cite{ducoffe-precioso:2018} mitigate
this by exploiting training dynamics with decision boundary based methodologies, 
but with a high cost at runtimes that makes them difficult at scale.


Lastly, compared to the dataset distillation, coreset selection is applied to a 
various domains.
It has been proven valuable for imbalanced datasets~\cite{sivasubramanian-aagalapatti-etal:2024,luo-zhao-etal:2024}, continual learning approaches~\cite{titsias-schwarz-etal:2019,mirzadeh-farajtabar-etal:2020}, 
machine unlearning~\cite{bourtoule-chandrasekaran-etal:2021} or even large languare models~\cite{killamsetty-zhao-etal:2021} and self-supervised approaches~\cite{kim-bae-yun:2023,singh-vashishtha-etal:2025}, further emphasizing the importance of effective and efficienct coreset selection methodologies.



% In this work, we take a different stance and investigate whether sustainable AI can be advanced through coreset selection rather than dataset synthesis. 
We argue that it is feasible to capture specific instances, within a training dataset, that are the most informative ones and preserve the networks decision information, with a lightweight approach, resulting faster and thus more energy efficient training.
We analyze the trade-off between subset size, accuracy, and training cost and show that carefully selected coreset can rival large-scale datasets in termes of accuracy. 


In summary the contributions of this works are:
\begin{itemize}
	\item a latent-space discrimination framework that leverages a pretrained diffusion model as a frozen encoder,
	\item a lightweight weigth embedder with a hyperplane assigned softmax head, ensuring class discriminative feature vectors,
	\item a cosine fartherst-first coreset selection approach for the discriminative feature space,
	\item an extensive experimentation on multiple networks, highlight the computational costs of different approaches for the coreset selection.
\end{itemize}


Data distillation, data selection and, in fact, most modern ANN
approaches follow a general pattern where the raw data is first mapped
to a learned latent representation that reduces the dimensionality of
the input. This approach offers the great benefit of mapping the data
to a more compact space where the transformations
applied by subsequent network layers can be more effectively learned.
Furthermore, since such latent representations are typically learned
in an unsupervised manner they are re-usable across learning tasks
over the same kind of data.

The work described here is positioned as follows with respect to the main
body of literature in dataset reduction:
%
\begin{itemize}

\item Data distilation: There should be no need to synthesise
  characteristic instances in a large dataset, we hypothesise that
  such a dataset already has the characteristic instances we need if
  we can select them,
\item Core set selection: Aiming for the most characteristic instances
  of each class will achieve the bulk of the accuracy, but is likely
  to miss nuanced distinctions. We hypothesise that selection should
  be driven by the difficulty or complexity of the distinctions that
  need to be learned, or, in other words, that selection should be
  biased not towards the most characteristic, central instances but
  towards the `difficult' cases, the ones that define the borders
  between classes.
\end{itemize}
%
Putting these two together, we hypothesise that a selection method can
be as effective as distilation (but at a lower computational cost),
and that the best strategy is to make a selection that is denser at
the periphery of a class than at its core.

